% Lab 7: Lab test. Compile: latexmk -xelatex lab07.tex
\documentclass[10pt,aspectratio=169,t]{beamer}
\usetheme[lab]{upbminimal}

\course{Application Development for Mobile Devices}
\decklabel{Lab 7}
\title{Lab test}
\subtitle{Rules, format, and two practice specs}
\author{Dinu-Ștefan Rusu}
\institute{FILS, Universitatea Politehnica București}
\date{Week 14}

\begin{document}

\titleframe

\begin{frame}[eyebrow=Today]{What this lab is for}
  \vskip 2mm
  \begin{grid}[2]
    \cell[\faIcon{stopwatch}]{Prove four skills, once}{Build a small app from a one-page spec in ninety minutes: a screen, a Cubit, a network call, and a route.}
    \cell[\faIcon{graduation-cap}]{Worth four points}{One point per requirement, graded from your pushed branch. There is no live demo.}
  \end{grid}
  \vskip 6mm
  \taskmeta{Time}{90 minutes}{Submit}{Push to admd-labs, branch lab-7, before time is up}
\end{frame}

\begin{frame}[eyebrow=Rules]{Individual, timed, your own machine}
  \begin{steps}
    \item You work alone. No pair programming, no shared screens, no copying a classmate's branch.
    \item Ninety minutes, one sitting, starting when you receive the spec.
    \item Your own laptop, your existing \code{admd-labs} repository from Labs 1 to 6.
    \item No new packages beyond what Labs 2 to 6 already added, unless the spec names one.
  \end{steps}
  \begin{takeaway}{The clock does not pause.}
    A crashed emulator or a broken hot reload eats into your ninety minutes. Restart calmly and keep going.
  \end{takeaway}
\end{frame}

\begin{frame}[eyebrow=Rules]{What you may and may not use}
  \begin{steps}
    \item Internet access is allowed: official docs, package pages, your own past labs.
    \item AI agents are allowed for writing code.
    \item You must be able to explain any line the grader points to, in your own words.
    \item Asking another person, in the room or online, to write code for you is not allowed.
  \end{steps}
  \begin{warning}
    An agent-written line you cannot explain earns no credit for that requirement, even if it compiles.
  \end{warning}
\end{frame}

\begin{frame}[eyebrow=Format]{One spec, four requirements}
  \vskip 2mm
  \begin{grid}[2]
    \cell[\faIcon{file-alt}]{One page, handed at the start}{An app idea and four requirements, on paper or a shared document. Nothing changes once the clock starts.}
    \cell[\faClock]{Ninety minutes, one push}{Build, commit as you go, and push to branch lab-7 before time is up. No extensions.}
  \end{grid}
  \vskip 6mm
  \muted{Every spec asks for the same four kinds of requirement, described next.}
\end{frame}

\begin{frame}[eyebrow=Format]{The four requirements, generically}
  \begin{grid}[2]
    \cell[\faThLarge]{A screen}{Built from the layout widgets: Row, Column, ListView, Card. No hardcoded pixel offsets.}
    \cell[\faToggleOn]{State in a Cubit}{An Equatable state class, \code{emit}, and a BlocBuilder. Not a raw setState.}
    \cell[\faWifi]{One network call}{A single HTTP call parsed into a model class with fromJson. Not a hardcoded map.}
    \cell[\faIcon{route}]{One navigation}{A go\_router route that carries data and pushes to a second screen.}
  \end{grid}
\end{frame}

\begin{frame}[eyebrow=Rubric]{One point per requirement}
  {\usebeamerfont{small}\begin{tabular}{@{}p{24mm}p{42mm}p{42mm}@{}}
    \toprule
    \thead{Requirement} & \thead{Full credit (1 pt)} & \thead{Partial credit (0.5 pt)} \\
    \midrule
    Screen & Matches the spec's widgets, no overflow & Builds, but the layout differs from spec \\
    Cubit state & UI updates through emit, never setState & Cubit exists, but a screen still calls setState \\
    Network + model & Parses the real response into a model & Call runs, but fields come from a raw map \\
    go\_router nav & Navigates and returns with the right data & Route exists, but push or data is wrong \\
    \bottomrule
  \end{tabular}}
  \vskip 2mm
  \muted{Anything not attempted scores zero.}
\end{frame}

\begin{frame}[fragile,eyebrow=Submission]{Push before the clock runs out}
\begin{shell}
git checkout -b lab-7
git push -u origin lab-7
\end{shell}
  \begin{steps}
    \item Create the branch the moment you receive the spec.
    \item Commit as you finish each requirement, with a message that names it.
    \item Push before the ninety minutes end.
    \item Open one pull request titled with your name. It does not need a review to count.
  \end{steps}
  \begin{warning}
    The grader reads the branch as it stands at the deadline. A late push is not graded.
  \end{warning}
\end{frame}

\begin{frame}[eyebrow=Prepare]{How to prepare}
  \vskip 2mm
  {\usebeamerfont{small}\begin{tabular}{@{}p{60mm}p{20mm}@{}}
    \toprule
    \thead{Requirement} & \thead{Taught in} \\
    \midrule
    Screen from the layout widgets & Lab 3 \\
    State in a Cubit & Lab 5 \\
    Network call and a model class & Lab 4 \\
    Navigation with go\_router & Lab 6 \\
    \bottomrule
  \end{tabular}}
  \vskip 4mm
  \begin{takeaway}{Nothing on the test is new.}
    Re-read the Task slides and the reference code in those four labs. Redo one small screen from memory before you come.
  \end{takeaway}
\end{frame}

\begin{frame}[eyebrow=Practice A]{Practice A: habit tracker}
  \muted{An app that lists daily habits the user can check off.}
  \vskip 2mm
  \begin{columns}[T]
    \begin{column}{0.56\textwidth}
      \begin{steps}
        \item A \code{ListView.builder} of habit cards, each with a checkbox, and a button that opens a dialog to add one.
        \item A \code{HabitsCubit} with an Equatable state holding the habit list. Toggling or adding emits a new state.
        \item On start, fetch \code{/users/1} from JSONPlaceholder and parse it into a \code{User} model shown in the AppBar.
        \item Routes \code{/} and \code{/habit/:id}. Tapping a card pushes to the detail route.
      \end{steps}
    \end{column}
    \begin{column}{0.4\textwidth}
      \kv{Grader checks}{}
      \begin{checklist}
        \item The app builds and runs.
        \item A checkbox toggle goes through the Cubit, never setState.
        \item The AppBar name comes from the API, not a hardcoded string.
        \item The back button returns from the detail screen to the list.
      \end{checklist}
    \end{column}
  \end{columns}
\end{frame}

\begin{frame}[eyebrow=Practice B]{Practice B: campus events board}
  \muted{An app that lists campus events, with filtering and a detail view.}
  \vskip 2mm
  \begin{columns}[T]
    \begin{column}{0.56\textwidth}
      \begin{steps}
        \item A row of filter chips (All, Today) above a scrollable list of event cards.
        \item An \code{EventsCubit} with a state holding the event list and the selected filter. Picking a chip emits a filtered state.
        \item On start, fetch \code{/posts} from JSONPlaceholder and parse each item into an \code{Event} model.
        \item Routes \code{/} and \code{/event/:id}. Tapping a card pushes to a detail screen with the full text.
      \end{steps}
    \end{column}
    \begin{column}{0.4\textwidth}
      \kv{Grader checks}{}
      \begin{checklist}
        \item A filter chip changes the list through the Cubit, not a local variable.
        \item Every field on screen comes from the Event model, never a raw map key.
        \item The detail route reads the id from the URL.
        \item The app does not crash when the network call is slow.
      \end{checklist}
    \end{column}
  \end{columns}
\end{frame}

\begin{frame}[eyebrow=Troubleshooting]{Mistakes that cost time}
  \vskip 2mm
  \trouble{Bad state: field does not exist for the class}{The JSON keys do not match the model. Print the raw response before writing fromJson.}
  \trouble{No routes matched location "/habit/3"}{The path pushed is not in the go\_router route list, or a name is misspelled.}
  \trouble{A blank screen, nothing in the console}{The Cubit never emitted after the network call finished. Check that the future actually completes.}
  \trouble{setState() or emit() called after dispose()}{A response arrived after the screen closed. Check \code{mounted} before calling either.}
\end{frame}

\closingframe{Push before time runs out.}{Branch lab-7 · one pull request · nothing after the deadline}

\end{document}
